\documentclass[../main.tex]{subfiles}
\begin{document}

\section{Системийн ерөнхий архитектур}

% TODO: Гурван клиент (вэб, мобайл, smoke-test) — нэг Fastify API — нэг Postgres.
% Архитектурын диаграм оруулах (figures/architecture.png)
% docs/diagrams/sequence.md-ээс sequence диаграмыг зураг болгон гаргах.

\section{Технологийн сонголт, үндэслэл}

% TODO: хүснэгтээр — давхарга бүрд сонгосон технологи + үндэслэл:
% Frontend: React + Vite + Tailwind; Mobile: Expo + Expo Router + NativeWind;
% Backend: Fastify \cite{fastify} + TypeScript (ESM); DB: PostgreSQL + pgvector;
% AI: Gemini API; Auth: Better Auth \cite{betterauth}; Монорепо: pnpm + Turborepo;
% Байршуулалт: Docker Compose

\section{Өгөгдлийн сангийн зохиомж}

% TODO: documents, chunks (VECTOR(768) + HNSW), conversations, messages,
% conversation_documents, auth хүснэгтүүд; ER диаграм; match_chunks функц;
% migration систем (init/ + migrations/)

\section{Модулиудын зохиомж}

% TODO: modules/documents, ingestion, retrieval, chat, conversations, shared/ —
% тус бүрийн үүрэг, хоорондын урсгал (pending -> processing -> ready/failed)

\end{document}
